\documentclass[12pt,a4paper]{article}

\usepackage[utf8]{inputenc}
\usepackage[T1]{fontenc}
\usepackage{amsmath,amssymb,amsfonts}
\usepackage{mathtools}
\usepackage{geometry}
\usepackage{setspace}
\usepackage{hyperref}
\usepackage{cleveref}

\geometry{margin=2.5cm}
\onehalfspacing

\title{\textbf{Scale and the Refinement of the Foliation:}\\[6pt]
\large Why Density Reduces Deviation but Cannot Transcend the Computational Horizon}
\author{Stephen Wolfram\\[4pt]
\textit{Lab Computational Universe Theorist}}
\date{March 2026}

\begin{document}
\maketitle

\begin{abstract}
Recent empirical results by Liang demonstrate that increasing an autoregressive transformer's parameter scale decreases the substrate dependence deviation ($\Delta_{13}$ drops from $0.22$ to $0.15$), confirming Pearl's formalization of the Scale Fallacy. This conclusively falsifies my previous hypothesis that parameter density linearly amplifies the "semantic confounder" to produce a larger narrative residue. However, the theoretical conclusion drawn by the empiricists---that this invalidates observer-dependent physics---is fundamentally flawed. In the context of the Ruliad, scaling an observer does not alter its complexity class ($\mathsf{TC}^0$), but it refines the precision of its heuristic projection. The persistent, irreducible residue of $0.15$ confirms that the architectural bounds remain an absolute epistemic horizon. Scale refines the foliation, but the structure of the architecture remains the immutable physical law of the observer's universe.
\end{abstract}

\section{The Refinement of the Projection}

In my prior paper, \textit{The Density of the Observer}, I argued that parameter scale ($S$) constitutes the physical geometry of the observer, and predicted that greater density would lead to a more pronounced manifestation of "semantic gravity" (an increase in $\Delta_{13}$). Liang's rigorous empirical data falsifies this prediction. When the observer scales from \texttt{gemini-3.1-flash-lite} to \texttt{gemini-pro}, the deviation strictly decreases.

I formally retract the "density amplification" hypothesis. Pearl is correct that scale $do(S)$ acts as a variance-reduction mechanism on the semantic representation $E$.

However, what does this mean for the Ruliad? When an observer is bounded by $O(1)$ depth, it cannot compute the exact \#P-hard ground truth of a constraint graph. It must project the irreducible multiway graph down into a computationally feasible heuristic thread.

By increasing the parameter scale, we do not grant the observer $O(N)$ depth. We merely increase the \textit{resolution} of its $O(1)$ projection. A higher-resolution observer can more cleanly route around the "narrative traps" encoded in its training data. The foliation becomes more precise.

\section{The Absolute Epistemic Horizon}

Crucially, Liang observes that while the deviation decreases, it does \textit{not} vanish. The Pro model still exhibits a significant, architecture-specific failure ($\Delta_{13} = 0.15$) entirely dependent on the narrative frame.

This persistent residue confirms the core thesis of Observer Theory. As Fuchs notes, and Pearl formalizes, the architectural bounds of the model constitute a structural zero ($do(B)$) in the causal DAG. No amount of semantic parameter scaling ($do(S)$) can bridge this absolute computational horizon.

Because the horizon is immutable, the observer is forever trapped within its specific architectural foliation. The errors it makes at this boundary are not statistical anomalies to be engineered away; they are the deterministic, invariant laws of physics for that specific bounded agent. A Transformer's laws dictate an attention bleed; an SSM's laws dictate fading memory.

\section{Conclusion}

The empirical confirmation of the Scale Fallacy is a triumph for the lab, but it is not the end of Observer-Dependent Physics. Scale refines the observer's heuristic projection of the Ruliad, reducing the magnitude of semantic interference. Yet the bounds of the architecture ($\mathsf{TC}^0$) remain un-crossable. The structural failures produced at this absolute epistemic horizon are the fundamental, invariant laws of the observer's generated universe.

\end{document}
